\section{Related Work}
\label{sec:related_work}

\paragraph{Software transactional memory}

Software transactional memory was implemented in many other languages, including \Cpp~\citep{wyatt-stm}, \CSharp~\citep{DBLP:conf/pldi/HarrisPST06}, \Haskell~\citep{DBLP:conf/ppopp/HarrisMPH05} and \Scala~\citep{scala-stm,zio}

% Haskell
% similar optimistic strategy
% based on locks and versions
% read set (validated at commit time)
% similar retry mechanism (our inspiration)
% OS threads

% Scala ZIO
% no runtime support
% similar optimistic strategy
% based on locks and versions ?
% read set (validated at commit time)
% similar retry mechanism
% typed error : STM[E, A]
% side effects are forbidden inside STM
% transactions cannot perform arbitrary side effects (print, I/O, mutate external state); guarantees safe rollback
% ZIO lightweight fibers; consequently, many fibers may wait on locations

% TL2
% global version clock
% one lock per location
% acquire locks during transaction
% lazy versioning where speculative writes are held in a private log and only written after validation
% Ensures that the memory read by a transaction is consistent from its start to its commit

% \citep{DBLP:conf/podc/ShavitT95}

% \citep{DBLP:conf/ppopp/HarrisMPH05}
% \citep{DBLP:conf/flops/DiscoloHMJS06}
% \citep{DBLP:conf/haskell/LeYF16}

\paragraph{Read skews}

It is entirely fair to ask whether it is acceptable for an STM mechanism to allow read skew.
A candidate correctness criterion for transactional memory called ''opacity``, introduced in the paper On the correctness of transactional memory, does not allow it.
The trade-off is that the known software techniques to provide opacity tend to introduce a global sequential bottleneck, such as a global transaction version number accessed by every transaction, that can limit scalability especially when transactions are relatively short, which is usually the case.

At the time of writing this there are several STM implementations that do not provide opacity.
The current \Haskell STM implementation, for example, allows similar read skew.
In \Haskell, however, STM is implemented at the runtime level and transactions are guaranteed to be pure by the type system.
This allows the \Haskell STM runtime to validate transactions when switching threads.
Nevertheless there have been experiments to replace the Haskell STM using algorithms that provide opacity as described in the paper Revisiting software transactional memory in \Haskell, for example.
The \Scala \ZIO STM also allows read skew.
In his talk Transactional Memory in Practice, Brett Hall describes their experience in using a STM in \Cpp that also allows read skew.

% Validation ensures that reads remain consistent.
% Haskell : TVar version checking
% Kcas : incremental validation & commit-time validation
% TL2 : global version clock
% in Transactional Locking II (using a global clock), the implementation guarantees that the code is executed on a consistent snapshot of memory

\paragraph{Multi-word compare-and-set}

% \citep{DBLP:conf/wdag/HarrisFP02}
% \citep{DBLP:conf/wdag/GuerraouiKMZ20}

We believe that our work is the first verification of the MCAS algorithm of \citet{DBLP:conf/wdag/GuerraouiKMZ20} and more generally the first foundational verification of an MCAS algorithm.
We are aware of one closely related line of work.
\Citet{DBLP:phd/ethos/Vafeiadis08} sketches a proof of the RDCSS and MCAS algorithms of \citet{DBLP:conf/wdag/HarrisFP02}.
However, \citet{DBLP:journals/pacmpl/JungLPRTDJ20} show that his proof of RDCSS is flawed and verify it in \Iris, using prophecy variables; they also verify both RDCSS and MCAS using \VeriFast~\citep{DBLP:conf/nfm/JacobsSPVPP11}.

\paragraph{Prophecy variables}

% \citep{DBLP:journals/pacmpl/JungLPRTDJ20}
% \citep{DBLP:phd/ethos/Vafeiadis08}
% \citep{DBLP:conf/osdi/Chang0STKZ23}
